\documentclass[10pt]{article}
\usepackage[margin=20mm]{geometry}
\usepackage{amsmath,amssymb,mathtools,fontspec,unicode-math,xurl,hyperref}
\IfFontExistsTF{DejaVu Serif}{\setmainfont{DejaVu Serif}}{\setmainfont{Latin Modern Roman}}
\setmathfont{latinmodern-math.otf}
\setlength{\parindent}{0pt}\setlength{\parskip}{6pt}\sloppy
\title{Scientific technical review of
The local structure of the energy landscape in multiphase mean curvature flow:
Weak-strong uniqueness and stability of evolutions}
\author{AI-assisted technical review}
\date{27 October 2026}
\begin{document}
\maketitle
Abstract: The assigned author preprint presents a relative-entropy method for weak-strong uniqueness of planar multiphase mean curvature flow. This review identifies two Major problems: a scope inconsistency in the varifold extension and a curvature sign that obstructs the stated compatibility estimates used in the global calibration construction. Three Minor problems concern entropy bookkeeping, the hypotheses of a Fourier criterion, and a paired series. An admissible shrinking-circle example disproves the printed local implication in Lemma 22; the incompatible normal derivatives in Proposition 33 establish its consequential use in the proof chain. The integral positivity needed for the actual Read–Shockley profile admits a complete analytic repair. These findings do not disprove the principal weighted BV uniqueness theorem. The overall assessment is Major; scientific adjudication and human review remain pending.\par
\section*{Identity, version, and theorem scope}
The authors are Julian Fischer, Sebastian Hensel, Tim Laux, and Theresa M. Simon. The journal association is DOI https://doi.org/10.4171/JEMS/1577. The exact assigned input is the 104-page arXiv:2003.05478v2, uploaded 27 September 2021. It is an author preprint, not the publisher Version of Record and not asserted here to be the accepted version. The publisher landing page separately records online-first publication in 2025, a Published date of 19 February 2025, submission on 13 December 2022, and acceptance on 26 December 2023. No equivalence with that later publisher text has been verified. PDF and printed page numbers coincide.\par
Theorem 1 (p. 11) concerns planar finite regular networks up to the strong solution’s lifetime. Definitions 14–16 (pp. 19–21) require regular interfaces, triple junctions, Herring balance, curvature and higher compatibility conditions, and a smooth family preserving the network structure. Constants on a fixed interval \(T<T_{\rm strong}\) are not claimed uniform at a topology change. Definition 9 (pp. 15–16) requires a symmetric tension matrix with positive off-diagonal entries, zero diagonal, and a nondegenerate Euclidean simplex realization \(\sigma_{ij}=|q_i-q_j|\). Strict triangle inequalities alone do not replace this requirement.\par
The BV competitors in Definition 13 (pp. 17–18) satisfy a transport identity, square-integrable normal velocities, a weighted first-variation identity, and a sharp global energy-dissipation inequality. The conclusion is not uniqueness among arbitrary Brakke flows. The abstract calibration implication (Theorem 3) and bulk-error estimate (Proposition 5) are conditional statements in \(d\ge2\); Theorem 6 and Lemma 8 construct the needed objects for the specified planar networks.\par
Write \[\xi_{ij}=\frac{\xi_i-\xi_j}{\sigma_{ij}},\qquad
 E_\sigma=\sum_{i\ne j}\sigma_{ij}\int_{I_{ij}}(1-n_{ij}\cdot\xi_{ij})\,d\mathcal H^{d-1}.\] The sums are over ordered pairs. The resulting double counting is intentional. Calibration coercivity gives \[|n-\xi|^2\le2(1-n\cdot\xi),\qquad
 0\le1-|\xi|^2\le2(1-n\cdot\xi),\] and controls the squared distance to the strong interface. Vanishing interface error gives containment; the signed transported weights in Definition 4 and their bulk error are also needed to identify the phases.\par
\section*{Established problems}
Problem 1 (Major: the varifold extension needs an explicit tension restriction).  Location. Definition 2(4a) and (5), pp. 12–13; Definitions 16 and 18, pp. 21, 23; (26), its following decomposition, and Theorem 19, p. 24; proof of Theorem 19, pp. 44–46, especially Step 4.\par
Evidence. Theorem 19 invokes the preceding strong solution and calibration without explicitly restricting the admissible \(\sigma_{ij}\) to unity. Definition 18 uses an unweighted integral-varifold mass \(\mu\). Its error functional is \[E_{\rm var}=2\mu(\mathbb R^d)-\sum_{i\ne j}\int_{I_{ij}}n_{ij}\cdot\xi_{ij}\,d\mathcal H^{d-1}.\] For \(\rho_i=d|D\chi_i|/d\mu\), the exact decomposition is \[E_{\rm var}
 =2\int\left(1-\frac12\sum_i\rho_i\right)d\mu
   +E_1,\qquad
 E_1=\sum_{i\ne j}\int_{I_{ij}}(1-n_{ij}\cdot\xi_{ij})\,d\mathcal H^{d-1}.       \tag{2.1}\] The text identifies the last term with the previously defined \(E_\sigma\). This is not an identity for general tensions. For example, take two phases, \(\sigma_{12}=2\), a unit-multiplicity smooth shrinking circle, and a different concentric reference shrinking circle over a common pre-extinction interval. The excess-multiplicity term is zero, and coercivity makes \(E_1>0\); nevertheless \(E_\sigma=2E_1\).\par
There is a substantive proof consequence beyond the displayed decomposition. Step 4 differentiates the unweighted cross term using \(-2\sum_i\int V_i\operatorname{div}\xi_i\,d|D\chi_i|\). This phase cancellation requires \(\xi_{ij}=\xi_i-\xi_j\), whereas Definition 2 uses the normalized difference above. Unequal tensions cannot in general be removed by changing all phase potentials. As exact local algebra, the admissible triangle with tensions \((1,1,\sqrt2)\) permits \[n_{12}=e_1,\quad n_{23}=e_2,\quad
 n_{31}=-(e_1+e_2)/\sqrt2.\] The weighted cyclic sum vanishes, but the unweighted sum is \((1-1/\sqrt2)(e_1+e_2)\ne0\). Thus these normals cannot be the three unweighted cyclic differences of phase vectors.\par
Impact and repair boundary. The written varifold extension is not justified under the preceding general-tension scope. State a unit-tension restriction (or an explicitly normalized equal-tension version), or introduce a suitable weighted varifold formulation and prove its dissipation and transport identities. A common constant tension can be normalized, so this is not evidence against the equal-tension result. The triangle calculation is a local cancellation obstruction, not a full counterexample to every hypothesis and conclusion of Theorem 19. No counterexample to the principal weighted BV theorem is asserted. The Major assessment concerns the material scope and proof input of this extension, not irreparability.\par
Problem 2 (Minor: two omissions in the entropy absorption step).  Location. Proposition 17(24), p. 22; definition of term II, p. 39; first display and (48), p. 40; Theorem 3(6), p. 13.\par
Evidence. Put \(w=V+\operatorname{div}\xi\). Replacing \(n\) by \(n-\xi\) in the mixed term involving \(I-\xi\otimes\xi\), as an equality, drops \[((I-\xi\otimes\xi)B)\cdot wn
 =((I-\xi\otimes\xi)B)\cdot w(n-\xi)
   +(1-|\xi|^2)(B\cdot\xi)w.                                  \tag{2.2}\] The last term need not vanish because a coercive calibration is short off its reference interface. At the purely algebraic values \(\xi=e_1/2,n=B=e_1,w=1\), the two terms on the right are \(3/8\) and \(3/8\), while the left is \(3/4\). This tests the printed identity, not a global flow.\par
The same page’s (48) also drops \(E_\sigma(0)\), present in Proposition 17. For distinct smooth concentric shrinking circles, \(E_\sigma(0)>0\) and \(E_\sigma(t)\) is continuous at zero. An estimate \(E_\sigma(T)+D(T)\le C\int_0^T E_\sigma(t)\,dt\), with nonnegative dissipation \(D\), cannot then hold as \(T\downarrow0\).\par
Repair and consequence. The existing assumptions suffice to repair both steps. With \(e=1-n\cdot\xi\) and bounded \(B\), \[|(1-|\xi|^2)(B\cdot\xi)w|
 \le\delta w^2+C_\delta(1-|\xi|^2)^2
 \le\delta w^2+2C_\delta e.\] Keep a sufficiently small dissipation fraction, restore \(E_\sigma(0)\), and apply Gronwall. This recovers the intended stability form; the omissions are not a counterexample to that corrected statement.\par
Problem 3 (Minor: the general Fourier criterion omits pointwise assumptions).  Location. Definition 9, pp. 15–16; (12)–(14) and Lemma 11, pp. 16–17; Fejér-limit argument, pp. 98–99.\par
Evidence. The generic part of Lemma 11 assumes a bounded profile \(f:[0,\pi/4]\to[0,1]\) and strict negativity of the nonzero Fourier coefficients of its even periodic squared extension. The proof later uses continuity of \(g=f^2\), absent from that statement. Define the continuous \(\pi/2\)-periodic function \[g_0(\theta)=\frac18\sum_{k\ge1}2^{-k}(1-\cos4k\theta),\qquad f_0=\sqrt{g_0}.\] Then \(0\le g_0\le1/4\), \(f_0(0)=0\), and its nonzero complex Fourier coefficients are \(\widehat g_0(k)=-2^{-|k|}/16<0\). Change only \(f(\pi/8)\) to zero, keeping \(f=f_0\) elsewhere on the defining interval, and extend evenly and periodically. The Fourier coefficients are unchanged. Two distinct orientations \(0,\pi/8\) nevertheless give \(\sigma_{12}=0\), violating strict off-diagonal positivity. The example already has zero diagonal.\par
If (12) is also to define the diagonal, a normalization is independently necessary: the continuous profile \(\sqrt{1/4+g_0}\) has the same strictly negative nonzero coefficients but \(f(0)=1/2\).\par
Repair and consequence. Require a continuous representative and \(f(0)=0\), or specify the zero diagonal separately. Fourier information alone is an almost-everywhere property and cannot control arbitrarily reassigned point values. With the pointwise hypotheses restored, the Fejér approximation and Vandermonde argument yield strict conditional negativity and the intended simplex realization. The actual Read–Shockley profile is continuous with \(f(0)=0\); this finding does not invalidate that subclass or Theorem 1, which independently assumes admissibility.\par
Problem 4 (Minor: an incorrect paired series in the Read–Shockley check).  Location. Equation (296), the two subsequent series, the \(500/11\) threshold, and the final numerical-integral assertion, p. 100.\par
Evidence. For \[I(\alpha)=\int_0^1(\log^2x+\log x-1)\cos(\alpha x)\,dx\] the unpaired series printed first gives \[I(\alpha)=\frac{5\alpha^2}{27}-\frac{7\alpha^4}{750}+O(\alpha^6).\] In the paired line the even-index numerator is printed \(4(2k)^2-6(2k)\); it must be \(4(2k)^2+6(2k)\). The printed expression instead has the coefficient \(-1/750\) at \(\alpha^4\). The correct first-pair positivity threshold is \(\alpha^2<1250/63\), not \(500/11\). Merely reversing the nearby malformed inequality does not bridge the resulting proof gap.\par
Complete analytic repair. In fact \(I(\alpha)>0\) for every \(\alpha>0\). Set \(w(x)=\log^2x+\log x-1\). Its integral is zero, and it changes sign once, from positive to negative, at \(r_0=\exp(-(1+\sqrt5)/2)\). For \(0<\alpha\le\pi\), the function \(\cos(\alpha x)\) strictly decreases on \((0,1)\). Therefore \[I(\alpha)=\int_0^1w(x)\{\cos(\alpha x)-\cos(\alpha r_0)\}\,dx>0.\] For the remaining range define \[\operatorname{Si}(a)=\int_0^a\frac{\sin t}{t}\,dt,\qquad
 J(a)=\int_0^a\frac{\operatorname{Si}(t)}t\,dt.\] Integration by parts, or termwise integration followed by integration of the sine integral, gives \[\alpha I(\alpha)=2J(\alpha)-\operatorname{Si}(\alpha)-\sin\alpha.              \tag{2.3}\] Pairing successive half periods of \(\sin t/t\) gives \(0<\operatorname{Si}(a)\le\operatorname{Si}(\pi)\) for \(a>0\); hence \(J\) is increasing. The bound \(\sin t\ge2t/\pi\) on \([0,\pi/2]\) gives \(J(\pi/2)\ge1\) and \(\operatorname{Si}(t)\ge1\) on \([\pi/2,\pi]\). Also \(\operatorname{Si}(\pi)\le\pi/2+\log2\), by splitting its integral at \(\pi/2\). Consequently, for \(\alpha\ge\pi\), \[\alpha I(\alpha)\ge2(1+\log2)-(\pi/2+\log2)-1
 =1+\log2-\pi/2>0.                                             \tag{2.4}\] For example, \(\log2>2/3\) and \(\pi<22/7\) make the last lower bound greater than \(2/21\).\par
Impact. The printed series argument is incorrect, but the required positivity and hence this ingredient of the actual Read–Shockley admissibility conclusion survive. The repair is analytic; no numerical integration or author program was executed or relied upon.\par
Problem 5 (Major: the tangential extension sign obstructs core compatibility estimates).  Location. Lemma 22, pp. 47–49, especially (65), (67), (75)–(76); (101)–(105), pp. 54–55; Proposition 33, (186)–(188) and Step 2, pp. 69–70; Construction 38, p. 83; Lemma 41, (268)–(269), pp. 87–89; Lemma 42, (271), (277), pp. 90–91.\par
Evidence. Use the source conventions \(\tau=J^\top n\), \(H=-\operatorname{div}_{\rm tan}n\), with \(H,n,\tau\) extended constantly in the normal direction, and \(B=Hn+\gamma\tau\). On the interface, \[\partial_\tau B=(\partial_\tau H+\gamma H)n
                  +(\partial_\tau\gamma-H^2)\tau,\qquad
 \partial_nB=(\partial_n\gamma)\tau .\] It follows that \[\nabla B:(n\otimes Jn+Jn\otimes n)
 =-(\partial_\tau H+\gamma H+\partial_n\gamma).                 \tag{2.5}\] To vanish there, as (76) requires, the right condition is \(\partial_n\gamma=-\gamma H-\partial_\tau H\). The printed (75) instead has \(+\gamma H-\partial_\tau H\), leaving the residual \(-2\gamma H\).\par
This is witnessed by a flow satisfying the local hypotheses, not merely arbitrary jets. Take two phases with unit tension and a shrinking circle \(R(t)=\sqrt{R_0^2-2t}>0\) on a closed interval before extinction. With the interior as the first phase, \[s=|x|-R(t),\quad n=e_r,\quad \tau=-e_\theta,\quad
 H=-1/R(t),\quad \gamma=\exp(-s/R(t)).\] These are smooth on a fixed sufficiently small tubular neighborhood. On the interface, \(\partial_n\gamma=\gamma H\) and \(\partial_\tau H=0\), so printed (75) holds. The first absolute-value term of (76) is \(2/R(t)>0\), the other is zero, and the claimed distance bound has right-hand side zero.\par
Actual core dependence. The error also prevents the stated matching between the curve and triple-junction velocity fields. Near an endpoint of the curve in Proposition 33, the cutoff in (188) is one. The interface value is therefore \(\gamma_c=\alpha\), where \(\alpha\) is the normally constant extension in (102). The triple-junction ansatz (101), with (105), has \[B^p=Hn+(\alpha+\beta s)\tau,\qquad
 \partial_n B^p=(-\alpha H-\partial_\tau H)\tau .\] But Proposition 33 expressly requires \(\gamma_c\) to satisfy the printed (75), and extends it by a normal Taylor expansion to do so. Thus \[\partial_n B^c=(+\alpha H-\partial_\tau H)\tau,\qquad
 \partial_n(B^p-B^c)=-2\alpha H\tau.                         \tag{2.6}\] At any point of the relevant interface segment where \(\alpha H\ne0\), this contradicts the \(O(\operatorname{dist})\) gradient estimate (187), whose right-hand side is zero on the interface. Along a normal segment it also gives \[B^p-B^c=-2\alpha Hs\tau+O(s^2),\] contrary to the \(O(\operatorname{dist}^2)\) estimate (186). The obstruction is not removed by choosing the tangential component to vanish identically: (104) imposes \(\partial_\tau\alpha=H^2\), so \(\alpha H\) cannot vanish throughout an open curved segment. The regular-network hypotheses do not require these segments to be flat. Step 2 on p. 70 identifies tangential values through (188), but equality on the interface does not supply the incompatible normal derivatives.\par
These estimates are actual inputs to the main construction. Construction 38 chooses the curve tangential component from Proposition 33 when the curve joins two triple junctions. Lemma 41 uses (186) to obtain (268), and uses (187) and (186) for the derivative estimate (269). Lemma 42 then explicitly invokes (269) in (277) and uses that identity to prove the global approximate transport estimate (271), required by the calibration condition (4c). This is a core proof defect as written, and its severity is Major even though a plausible correction is short. The bound used in that argument and the conclusion it was intended to prove must nevertheless be distinguished: the leading gradient discrepancy is \(-2\alpha H\,\tau\otimes n\), whose transpose annihilates \(n\). Thus this discrepancy alone does not refute the projected transport condition (271); a weaker projected estimate might suffice in a repaired proof.\par
Repair and limit. Reverse the sign of \(\gamma H\) in (75) and use the corresponding normal Taylor coefficient in Proposition 33. This agrees with the already printed \(\beta=-\alpha H-\partial_\tau H\) in (105), a substantive consistency check supporting the intended construction. A complete verification of the corrected global gluing argument is outside this focused acceptance. The shrinking circle is an actual counterexample to the local implication; the normal-jet calculation establishes the specified compatibility obstruction. Neither is asserted to be a counterexample to the corrected construction or to the main uniqueness theorem.\par
\section*{Positive checks, dependencies, and limits}
The entropy controls, the ordered-pair factors, and the bulk-weight sign identity were checked. If \(\vartheta_i\) is negative inside the strong phase and positive outside, then \(|\chi_i-\bar\chi_i||\vartheta_i|=(\chi_i-\bar\chi_i)\vartheta_i\). The strong-interface transport term is killed by \(\vartheta_i=0\) there. The remaining velocity and distance terms admit the advertised Young estimates. Both interface and bulk errors must vanish to obtain phase equality almost everywhere.\par
The local signed-distance transport identity \((\partial_t+B\cdot\nabla)s=0\) yields \((\partial_t+B\cdot\nabla)n+(\nabla B)^\top n=0\). At a triple point with velocity \(p'\), the relations \(p'=H_i n_i+\alpha_i\tau_i\), \(\partial_{\tau_i}\alpha_i=H_i^2\), and (105) explain the matching first jets: \[\nabla\widetilde\xi_i=\tau_i\otimes Jp',
 \qquad
 \nabla\widetilde B_i
 =\beta_i(\tau_i\otimes n_i-n_i\otimes\tau_i)=-\beta_iJ .\] The higher compatibility condition supplies a common \(\beta_i\) at the junction. These checks support the construction’s mechanism; they are not a complete reconstruction of every interpolation and global gluing estimate. In particular the preliminary quadratic normal ansatz has squared length \(1+\alpha_i^4s_i^4/4\), so its later normalization cannot be silently omitted.\par
The transported weights need not be compactly supported to be integrable and strictly signed. The construction on pp. 95, 98 multiplies by \(\eta_R(x)=\kappa(e^{R-|x|})>0\), with \(|\nabla\eta_R|\le C\eta_R\). This resolves a potential concern about the exterior phase. The reduced-boundary/topological-boundary equality in Definition 14 deserves a closure or almost-everywhere convention at triple points; it is not counted as a separate theorem defect here.\par
\section*{Assessment and actual review coverage}
The assessment is Major, based on Problems 1 and 5. Problems 2–4 remain Minor with the explicit consequences and repairs established above. Problem 5 is Major because the incorrect extension enters the core calibration construction; the short form of its proposed correction does not reduce that impact. No claim is made that the main weighted BV uniqueness theorem is false. The exact reviewed version is older than the journal publication.\par
The initial-review principal checked the main definitions and theorem chain, entropy and bulk-error steps, the varifold cancellation, signed-distance and triple-point formulas, the transported-weight tail, and the Fourier/positivity argument. The initial review recorded actual inspection of thirty-six source evidence pages as rendered images: 1, 11–24, 39–46, 48–52, 54–55, 69, 94–95, and 98–100. The full 104-page proof was not independently reconstructed: substantial parts of Sections 3, 6.2, and 7, and the cited external existence/regularity results, remain dependencies rather than independently certified proofs. Text extraction and manifest parsing are not counted as scientific proof verification.\par
Their actual scopes and hashes are retained separately; they are not separately authenticated Pro reviewers or independent acceptance authorities. The external controller’s UI observation, this runtime’s self-description, and the absence of independent backend verification are recorded separately.\par
The initial review consulted the supplied same-version TeX and figures selectively. Public publisher/arXiv metadata and a bounded correction/supplement search were checked; no relevant correction or separate executable supplement was identified within that search. Absence of code is not a mathematical defect. No author code was run. The focused independent acceptance rechecked the five reported problems, their hypotheses, counterevidence, and consequences against the exact assigned source. It followed the additional dependency chain for Problem 5 through pp. 83 and 87–91, revised its severity and impact, and preserved the initial report. This is not a new complete review of all 104 pages. The current scientific status remains `scientific\_decision=pending` and `human\_reviews=[]`.\par
9 J. Fischer, S. Hensel, T. Laux, and T. M. Simon, The local structure of the energy landscape in multiphase mean curvature flow: Weak-strong uniqueness and stability of evolutions, arXiv:2003.05478v2 (27 September 2021), assigned 104-page PDF and supplied same-version TeX. arXiv version record, https://arxiv.org/abs/2003.05478v2, accessed 26 October 2026. EMS Press, publication record for DOI 10.4171/JEMS/1577, https://ems.press/journals/jems/articles/14298604, accessed 26 October 2026.\par
\end{document}
